\chapter{Results: Strings and Winds from Tracked Audio}
\label{ch:phase2}\label{sec:phase2}
This chapter is stage 2 of the observation ladder, in which the prior of
the piano setting meets realistic measurement uncertainty. It carries the
proposed model to strings and winds, where the performer controls pitch.
Nothing in the latent prior changes. The graph and feature kernels of
\eqref{eq:gpfirst}, its first two terms, the hyperparameter fit
(\S\ref{sec:eb}) and the exact conjugate posterior (\S\ref{sec:gpfirst})
are reused. One term is absent: the music model is trained on piano, so
its embedding kernel takes no part here, and the evidence branch consists
of the hand-built score features alone. The observation-noise term belongs
to the likelihood, and the likelihood is what changes. Three things
therefore change, and the first section specifies each of them: the
channel set, the way the targets are obtained, and the observation noise,
which stops being negligible.


\section{The model in this setting: channels, targets and noise}
\label{sec:phase2-delta}

\subsection{The channel set}
A fixed-pitch keyboard affords no performer-controlled intonation, and its
dynamics is a MIDI velocity. Neither holds for strings and winds. We call
the per-note vector of channels the \textbf{bundle}, in the everyday sense
fixed in the Introduction, and in this setting it becomes
\begin{equation}
  y_i \;=\; \bigl[\, \tau_i,\ \log r_i,\ \ell_i,\ c_i,\ f_i^{\mathrm{vib}},\ \gamma_i,\ d_i^{\mathrm{vib}} \,\bigr]
  \;\in\; \R^{k},
  \label{eq:phase2-channels}
\end{equation}
where each channel is, as on piano, one scalar per note of the known
score. The first two channels keep their piano meaning: $\tau_i$ is the
onset residual against the score-time warp, re-anchored for this corpus
by the annotated onsets as described below, and $\log r_i$ is the
articulation ratio. The loudness channel $\ell_i$ replaces the MIDI
velocity $v_i$ and is read from the note's own audio. The note's span is
cut into four equal-length chunks of samples $x^{(q)}$, $q = 1, \dots, 4$,
the energy of each chunk is summarised by its log root-mean-square, and
\begin{align}
  \ell_i &\;=\; \frac{1}{4} \sum_{q=1}^{4} \log \mathrm{rms}\bigl(x^{(q)}\bigr)
  \;-\; \overline{\ell}_{\mathrm{track}},
  \label{eq:loudness}\\
  \widehat{\operatorname{Var}}[\ell_i] &\;=\; \frac{s^2_{\mathrm{chunks}}}{4}.
  \nonumber
\end{align}

The chunking is what makes loudness an estimate rather than a bare
number. The four chunks are internal replicates, so their mean is the
note's level, and the sample variance of the four log chunk-RMS values,
$s^2_{\mathrm{chunks}}$, yields the standard error
$s^2_{\mathrm{chunks}}/4$. This standard error is automatically larger for
notes whose level changes inside the note, during the attack or the decay,
which is exactly when one loudness value per note is least well defined.
Four is an untuned small default, enough replicates to estimate a spread
while the chunks remain long enough to average over waveform periods on
short notes. The term $\overline{\ell}_{\mathrm{track}}$ is the track's
mean of the per-note values, and subtracting it centres the channel per
recording. Absolute level is arbitrary per track, through microphone gain,
distance and mastering, so only relative loudness within a recording is
meaningful. The remaining channels are new to this setting. The static
intonation offset $c_i$ measures how sharp or flat the note's held centre
sits, and $(f_i^{\mathrm{vib}}, \gamma_i, d_i^{\mathrm{vib}})$ are the
vibrato rate, extent and onset delay defined below. The superscript keeps
the delay distinct from the performed duration $\delta_i$ of
\eqref{eq:logr}.

This vector leaves two decisions open, and one of them is settled by
measurement against a criterion fixed before the numbers were computed.
The delay is identifiable on $95$--$97\%$ of vibrato-identifiable notes
on both the reference and the tracked pitch curves, with a median
agreement of $18$\,ms between tracker and reference, so
$d_i^{\mathrm{vib}}$ is a channel and $k = 7$ as written. On the
evaluation corpus (\S\ref{sec:urmp}) articulation has no measured offset
chain, that is, no sequence of annotated performed note ends from which
$\log r_i$ could be computed, so the bundle evaluated there consists of
the remaining six channels. The second decision stays open: whether
$\ell_i$, which is an amplitude rather than a keystroke, is comparable
enough with $v_i$ for the two to be discussed as one channel across
settings. We do not assume that they are.

What these quantities mean to a listener, how sharp or flat and by how
much, the vibrato's speed, width and delayed onset, the note's loudness
relative to its own recording, is spelled out in musical terms in
Appendix~\ref{app:glossary}. Structurally, what matters is that none of
these values is annotated by hand. Every one is \emph{estimated} from the
recording: a pitch tracker supplies the note's pitch curve, and the
per-note fit described next turns that curve into these numbers, each
with its own uncertainty. This is why the observation noise of
\eqref{eq:phase2-noise} becomes substantial in this setting. The corpus's
human annotations supply only the note spans, that is, the onsets,
durations and written pitches. Its annotated pitch curves serve to
calibrate the tracker and never as channel values.

\subsection{From a pitch curve to per-note targets}
Let $f_0^{(i)}(t)$ be the fundamental estimated on the sounding interval
of note $i$, restricted to the samples the tracker marks as voiced. Its
deviation from the equal-tempered target, in cents, is
\begin{equation}
  c_i(t) \;=\; 1200\,\log_2 \frac{f_0^{(i)}(t)}{f_{\mathrm{ref}}\, 2^{\,(p_i - 69)/12}} ,
  \label{eq:cents-curve}
\end{equation}
with $f_{\mathrm{ref}} = 440$\,Hz the frequency of A4, whose MIDI number
is $69$, so that the denominator is the equal-tempered frequency of the
written pitch $p_i$. This curve is not itself a target, because it mixes
a static mistuning with an oscillation. Following the standard sinusoidal
description of vibrato \cite{seashore1936,prame1994,prame1997}, it is
decomposed as
\begin{equation}
  c_i(t) \;=\; \underbrace{c_i}_{\text{intonation}} \;+\; u_i(t) \;+\; \epsilon_i(t),
  \qquad
  u_i(t) \;=\;
  \begin{cases}
    \gamma_i \sin\!\bigl(2\pi f_i^{\mathrm{vib}}\,(t - t_i - d_i^{\mathrm{vib}})\bigr), & t \ge t_i + d_i^{\mathrm{vib}},\\[2pt]
    0, & \text{otherwise},
  \end{cases}
  \label{eq:vibrato}
\end{equation}
with $t_i$ the performed onset. Thus $c_i$ (cents) is the vibrato-free
centre, $\gamma_i$ (cents) the vibrato extent, $f_i^{\mathrm{vib}}$ (Hz)
its rate, and $d_i^{\mathrm{vib}}$ (s) the delay before it begins.

Estimating $(c_i, f_i^{\mathrm{vib}}, \gamma_i, d_i^{\mathrm{vib}})$ is a
small nonlinear least-squares fit per note on the voiced samples, and
per-note vibrato parameter estimation from recordings is an established
analysis task \cite{driedger2016}. The fit runs in two forms. A continuous
sinusoid supplies $(c_i, \gamma_i, f_i^{\mathrm{vib}})$, and in that form
the delay coordinate is only a phase, identified modulo one period. The
delayed-onset form of \eqref{eq:vibrato}, in which the oscillation grows
out of the centre at $t_i + d_i^{\mathrm{vib}}$, supplies the delay
itself, with its own identifiability rule and a variance obtained by
linearisation. The centre $c_i$ is \emph{not} the mean of $c_i(t)$: the
mean coincides with the vibrato-free centre only over a whole number of
cycles, which a note rarely supplies, and perceived pitch likewise departs
from the curve mean on short vibrato tones \cite{dalessandro1994}.

The observation rule fixes every threshold by the development
measurements of \S\ref{sec:urmp}. Note $i$'s sample set is
\begin{equation}
  J_i \;=\; \bigl\{\, \text{frames } j:\; t_j \in [t_i,\, t_i + \delta_i),\;
  \text{voiced}_j,\; \pi_j \ge q_{0.2}(\text{track}) \,\bigr\},
  \label{eq:framerule}
\end{equation}
so that the tracker's per-frame confidence $\pi_j$ is thresholded at the
track's lowest quintile, which the development measurements identified as
the gross-error filter. From this set follow three observation decisions,
stated separately.
\begin{itemize}[leftmargin=\parindent,noitemsep]
  \item The intonation entry is observed if and only if $|J_i| \ge 4$ and
        $|\hat c_i| \le 150$ cents, the octave-failure rule validated
        against the reference.
  \item The vibrato entries additionally require identifiability, at
        least $8$ voiced frames and at least $1.5$ cycles at the fitted
        rate, without which $(\gamma_i, f_i^{\mathrm{vib}},
        d_i^{\mathrm{vib}})$ cannot be read from the note at all.
  \item Every other entry is missing for the model.
\end{itemize}
One choice remains open, namely a robust loss rather than least squares
inside the fit.

\subsection{Per-channel missingness}
The identifiability rule makes missingness structural rather than
incidental. A short note may support $c_i$ but no vibrato parameters, so
a target can be missing for some channels of a note and not for others.
We call each (note, channel) pair an \textbf{entry}, and the resulting
pattern of observed entries the \textbf{observation mask}. The conjugate
posterior of \S\ref{sec:gpfirst} assumes per-note masks, which is exactly
what keeps the observed block Kronecker-compatible (\S\ref{sec:gpfirst}).
Per-entry missingness breaks that structure but not the exactness: the
conjugate posterior conditions on an arbitrary observed index set, so
what must give way is the per-note mask interface and any
Kronecker-structured shortcut, not conjugacy. The model therefore accepts
per-entry observation masks. Its exactness in this generalised form is
verified against brute-force Gaussian conditioning, the per-note case
reproduces the published path, and a synthetic end-to-end pilot exercises
it with estimator-induced missingness (\S\ref{sec:phase2-pilot}).
Figure~\ref{fig:phase2-datapoint} shows one data point of each kind.

\figph{phase2_datapoint_dev.png}{1.0}{One strings-and-winds data point (violin,
development track). Top: a $1.49$\,s note, showing the confidence-kept
tracker frames (banded at the tracker's $\sim$10-cent frequency grid),
the delayed-onset fit of \eqref{eq:vibrato} riding them, the flat stretch
before the fitted onset delay (dashed) and the vibrato-free centre
(dotted). At right, the note's six-entry record with each entry's
estimator uncertainty, where the wide $\pm 37$\,ms on $\tau$ is the
warp's noise row. Bottom: a $0.32$\,s note, too short for the
identifiability rule, so that its three vibrato entries are missing while
intonation, loudness and timing are still observed, which illustrates the
observation mask. The third row shows the loudness construction of
\eqref{eq:loudness} on the same long note: the waveform quartered into
equal chunks, each chunk's $\pm$RMS level and log-RMS value, and the
computation, in which the chunk mean minus the track mean gives
$\ell_i = -0.18$ and the chunk spread gives the $\pm 0.12$ standard
error, widened by the audible crescendo inside the note. The result
matches the $\ell$ row of the record card.\label{fig:phase2-datapoint}}

The missingness mechanism deserves naming \cite{rubin1976}. The
evaluation's hidden notes are missing completely at random by
construction, through a seeded coin per note. The estimator's missing
entries are not: the identifiability rule removes a vibrato entry exactly
when the note's vibrato is short or weak, so this missingness is
informative, missing \emph{not} at random and dependent on the value it
hides. Two consequences follow. The held-out scores stay interpretable,
but for the identifiable sub-population only: random masking inside an
informatively filtered set of entries is unbiased for the notes whose
vibrato the estimator can measure, not for all notes. And the posteriors
the model supplies at estimator-missing entries borrow from vibrato-rich
observed neighbours and may therefore be biased toward audible vibrato,
of which the delay posterior's occasional sub-zero stretches are the
visible signature. Filled entries are posterior extrapolations under the
prior and are to be read as such, not as recovered measurements.

\subsection{Heteroscedastic observation noise}
On piano the targets are read from clean MIDI and $\Sigma_e \approx 0$.
Here every target is an estimate, and its uncertainty is both known and
unequal across notes: the $f_0$ tracker reports a per-frame periodicity
or confidence, and the per-note fit of \eqref{eq:vibrato} returns a
parameter covariance. Writing $\sigma_{e,ij}^2$ for the variance of
channel $j$ of note $i$,
\begin{equation}
  \tilde y \;=\; y + e, \qquad
  e \sim \N\bigl(0,\ \Sigma_e\bigr), \qquad
  \Sigma_e = \operatorname{diag}\bigl(\sigma_{e,ij}^2\bigr),
  \label{eq:phase2-noise}
\end{equation}
which replaces the per-channel term $\operatorname{diag}(\varsigma^2)\otimes I$ of
\eqref{eq:gpfirst} by a full diagonal over (note, channel). Conjugacy is
untouched, since a diagonal is a diagonal, so the posterior and the
marginal likelihood are unchanged in form. The noise floor and the
safeguard of \S\ref{sec:eb} carry over.

One decision remains: whether the estimator-supplied variances enter the
likelihood unchanged, or carry a per-channel learned scale, so that the
evidence may discover that the tracker is over- or under-confident. We
call the first option the \textbf{given variances}, the estimator's own
variances used without rescaling. The synthetic pilot
(\S\ref{sec:phase2-pilot}) reverses the a-priori guess. With calibrated
estimator variances, the given variances are better on every channel,
and the learned scale over-shrinks the vibrato-extent channel (coverage
$0.76$--$0.81$ against $0.86$--$0.87$). The learned scale is therefore
the fallback for a mis-calibrated tracker, not the default. Real data
gave the same answer (\S\ref{sec:phase2-real}): with estimator octave
failures present, the learned scale collapses while the given variances
stay healthy, and after the octave-failure rule the given variances
remain the better-calibrated choice.

\subsection{What the per-note compression discards}
The channels of \eqref{eq:phase2-channels} hold every parameter fixed
within the note: they are time-averages. A development drift study
(Appendix~\ref{app:drift}) measures what the averaging discards. About
two thirds of notes carry a real within-note drift of the intonation
centre, found by the tracked and the reference curves independently,
with a median of $10.5$ cents per note, an order of magnitude above the
entry's reported standard error. Yet the discarded structure is
graph-white: within-note slopes carry essentially no note-to-note
correlation ($+0.03$, against $+0.59$ for timing, the channel the graph
demonstrably helps). The per-note resolution is therefore empirically the
right level for this model, and the mismatch is priced into the
residual-based entry variances, which is why calibration held on the
held-out test set. The discarded structure itself belongs to the waveform
setting's frame-level likelihood, whose amplitude envelope $A_i(t)$ and
pitch-curve model are built to carry it. Slope channels remain candidates
for later work under the $d^{\mathrm{vib}}$ template, included as output
and asserting nothing.

\subsection{The score-time warp and the alignment error}
The piano setting obtained the score-time warp $T(\cdot)$ from an aligned
beat grid supplied with the MIDI. There is no MIDI here. $T$ must come
from audio-to-score alignment, and its error propagates directly into
$\tau_i$, since $\tau_i = t_i - T(b_i)$ by definition. Worse, alignment
error is correlated along score time, which is precisely the structure
the graph prior expects expressive timing to have. A model fit naively
will therefore absorb some alignment error as expressive timing and
report it with confident error bars. There are two defensible responses.
One folds the aligner's own uncertainty into the $\tau$ row of
\eqref{eq:phase2-noise}, with the caveat that \eqref{eq:phase2-noise} is
diagonal, so it captures the error's scale but cannot represent its
correlation along score time. The other withholds $\tau$ from this
chapter's hypotheses entirely and reports only the channels the audio
measures directly. This is the single largest threat to validity in this
setting, and it has no analogue on piano.

On the chosen corpus the threat has a measured resolution. URMP's
note-level onset annotations anchor the warp without any audio aligner on
76 of 78 development tracks, half by exact score-order matching and half
by a pitch DTW. The two tracks without a usable warp remain in the
evaluation with their $\tau$ entries missing, not dropped. The
feasibility statistic uses a local leave-one-out warp: with score-onset
and performed-onset pairs $(b_j, t_j)$, each note gets a tempo line with
slope $\hat\varphi_i$ and intercept $\hat\psi_i$ fitted from its
neighbours alone,
\begin{align}
  (\hat\varphi_i,\, \hat\psi_i)
  &\;=\; \arg\min_{\varphi,\,\psi}\; \sum_{0 < |j - i| \le 8}
     \bigl(t_j - \varphi\, b_j - \psi\bigr)^2,
  \label{eq:localwarp}\\
  \tau_i &\;=\; t_i \,-\, \bigl(\hat\varphi_i\, b_i + \hat\psi_i\bigr),
  \nonumber
\end{align}
so the note being judged never enters its own warp. The result is a
median $\operatorname{std}(\tau)$ of $79$\,ms with lag-1 neighbour
correlation $+0.59$, that is, piano-scale timing with exactly the
structure the graph prior models. On this basis $\tau$ is a channel of
the evaluated bundle, and its noise row is the tempo line's leave-one-out
predictive variance, which carries the aligner error's scale but not its
correlation, as noted above. This is development evidence.

\subsection{The score graph and the tonal metric}
The graph construction \eqref{eq:adjacency} and its alternatives
(\S\ref{sec:graphvariants}) carry over, with per-channel length scales.
One earlier result becomes newly interesting. The kernel comparison
(\S\ref{sec:kernels}) found that replacing the semitone metric by the
circle-of-fifths metric \eqref{eq:tonaladj} made recovery and calibration
significantly worse: expressive behaviour travels with register
proximity, not tonal proximity. But that experiment measured timing,
articulation and velocity, none of which is a pitch. This setting
introduces a channel whose target is pitch itself. It is therefore an
open and cheap question whether the tonal metric that is worse for
expression is better for intonation, that is, whether notes a fifth apart
share a mistuning that notes a semitone apart do not, as temperament and
a fixed tuning reference would predict.

The test was run on the development set, with the hypothesis stated
before the numbers were computed (Appendix~\ref{app:tonal}). The
circle-of-fifths metric improves held-out intonation significantly on
both axes while re-imposing the familiar replacement penalty on timing.
The tonal geometry helps exactly where the target is pitch and hurts
where it is not, the first evidence in this thesis that a music-theoretic
\emph{geometry}, rather than merely music-theoretic \emph{edges}, earns
its place. The hypothesis was then evaluated on the held-out test set
(\S\ref{sec:phase2-reuse}): the recovery gain held, but the calibration
gain, the pre-specified primary measure, did not, its interval reaching
zero. The metric stays exploratory, and none of the main hypotheses of
this chapter depends on it.


\section{Data and protocol}\label{sec:phase2-protocol}
The corpus is URMP, described with the other corpora in \S\ref{sec:urmp},
together with the three preparatory measurements that govern its use:
tracker calibration, the composition-level split into a development set
and a held-out test set, and the sense in which recovery means agreement
with the estimator's targets rather than with a gold standard. That sense
applies to every number in this chapter. The evaluation protocol is that
of the piano setting, held-out imputation on both axes with a paired
per-piece bootstrap (\S\ref{ch:data}), and the selection discipline of
\S\ref{sec:hygiene} applies without exception. Because the targets carry
real noise, calibration rather than RMSE becomes the more informative
axis, and the predictive variance floor is no longer a technicality but a
first-order modelling term. The held-out test set of this setting is the 13-piece set of
\S\ref{sec:hygiene}, and the hypotheses it tests
(\S\ref{sec:phase2-real}) were fixed on the development set alone.


\section{Development results}\label{sec:phase2-devres}

\subsection{A synthetic pilot}
\label{sec:phase2-pilot}
Before any real audio, the full pipeline was validated end to end on
synthetic pieces whose per-note truth is known. On these pieces the graph
model is better than the no-graph ablation on every channel. Missing
vibrato entries, which no per-note method can even express, are recovered
at roughly $4\times$ better than prior scale with coverage at nominal,
and the given-variance noise default is settled there. This is a
synthetic result: it validates the machinery and settles design
defaults, and it asserts nothing about real audio.

\subsection{Results on real audio}
\label{sec:phase2-real}
The full pipeline then ran end to end on real recordings: tracked $f_0$,
confidence-filtered per-note estimation of $[c, \log\gamma, \log
f^{\mathrm{vib}}]$, per-note loudness $\ell$ (log RMS, centred per
recording), the onset-anchored timing channel $\tau$, and the delay
$d^{\mathrm{vib}}$ from the delayed-onset fit, which together form the
six-channel bundle. The estimator variances enter as the given variances,
with per-entry missingness from the identifiability rule and the
heteroscedastic model with per-entry observation masks. The per-track
vibrato-identifiable fraction ranges from $2\%$ to $97\%$ with median
$38\%$, so missingness is the norm, not an edge case. The run covers 77
of the 78 unique development tracks, after deduplication of shared
recordings and with one track below the minimum of 30 usable notes,
hiding 30\% of notes. One target-quality rule was added with measured
justification rather than by assumption: $2.2\%$ of intonation targets
are octave slips ($|c| \approx 1200$ cents), the cross-check against the
reference rejects every one of them (0 of 358 agree with the
reference-derived value), and they are marked missing, treated exactly
like unidentifiable vibrato.

Paired against the no-graph ablation under the given-variance noise
model, the declared default, recovery improves significantly on
intonation and on both vibrato channels ($\Delta$RMSE $-0.89$ cents,
$-0.26$, $-0.30$), calibration improves significantly on both vibrato
channels and on timing ($\Delta$NLL $-3.8$, $-0.43$, $-0.29$), and
coverage is $0.88$--$0.91$ at nominal $90\%$ on every one of the six
channels. Table~\ref{tab:phase2-dev} collects the paired contrasts. Under
this noise model the loudness, timing and delay recovery contrasts are
small and not significant ($-0.003$, $-0.002$, $-0.001$), and the
loudness NLL contrast is $+0.04$ \emph{against} the graph, small but the
bundle's one significant adverse result. In the per-family breakdown,
computed under the learned-scale contrast, the brass intonation contrast
is not significant on its own, while the vibrato gain holds within every
family and the intonation gain within strings and
woodwind.\footnote{The learned-scale fallback makes all six recovery
contrasts significant but is the worse-calibrated noise model against the
reference (coverage $0.72$--$0.86$ vs $0.82$--$0.86$). Its full rows are
in the results record.} The ordering of systems is unchanged when scored
against the reference targets derived from the provided pitch curves.

A check with fresh mask seeds (seeds 2, 3, same protocol) reproduces
every direction, every significant recovery difference to two decimals,
and the adverse loudness result. The significance of the heavy-tailed
calibration contrasts changes with the seed: the extent-channel
$\Delta$NLL, $-3.8$ with an interval excluding zero on the first seed, is
$-0.22$ and not significant on the fresh seeds. This variability was
therefore known before the held-out evaluation rather than discovered by
it.

\begin{table}[t]\centering
\caption{Strings-and-winds development results under the given-variance
noise model, the declared default, paired per (track, seed) against the
no-graph ablation ($95\%$ bootstrap intervals, with $^{*}$ marking an
interval that excludes zero and $n = 148$--$153$ pairs per channel), with
coverage at nominal $90\%$. All numbers are development numbers. The
hypotheses of the held-out evaluation were fixed on these columns. The
fallback noise model is reported in the results record.}
\label{tab:phase2-dev}
\footnotesize
\begin{tabular}{@{}lrrc@{}}
\toprule
channel & $\Delta$RMSE & $\Delta$NLL & cov@90 \\
\midrule
intonation $c$ (cents) & $-0.891^{*}$ & $-0.020$\phantom{$^{*}$} & 0.89 \\
extent $\log\gamma$ & $-0.256^{*}$ & $-3.772^{*}$ & 0.90 \\
rate $\log f^{\mathrm{vib}}$ & $-0.300^{*}$ & $-0.426^{*}$ & 0.88 \\
loudness $\ell$ & $-0.003$\phantom{$^{*}$} & $+0.042^{*}$ & 0.90 \\
timing $\tau$ (s) & $-0.002$\phantom{$^{*}$} & $-0.292^{*}$ & 0.91 \\
delay $d^{\mathrm{vib}}$ (s) & $-0.001$\phantom{$^{*}$} & $+0.057$\phantom{$^{*}$} & 0.89 \\
\bottomrule
\end{tabular}
\end{table}

Estimator-missing vibrato entries on visible notes receive posteriors,
the capability no per-note method has (Fig.~\ref{fig:phase2-real}). All
of this is development evidence in the piano chapter's sense. It selected
the design defaults above and fixed the hypotheses for the held-out
evaluation. The main hypotheses are recovery on intonation, calibration
on \emph{both} vibrato channels, and coverage within $[0.85, 0.95]$ on
all six channels. Timing calibration is a secondary hypothesis whose
outcome does not decide the others, and the delay channel carries no
hypothesis. Every number of the held-out evaluation is reported.

\section{Held-out evaluation}
\label{sec:phase2-conf}
The held-out test set, 13 pieces and 40 unique tracks after recording
deduplication, was evaluated once with two mask seeds under the protocol
of \S\ref{sec:hygiene}, and Table~\ref{tab:phase2-conf} reports the
outcome. The main hypotheses were supported. Intonation recovery is
$-0.877\,[-1.236, -0.538]$, an interval that excludes zero, reproducing
the development basis of $-0.891$ almost exactly. Vibrato calibration
improves on both channels, $-2.990\,[-6.575, -0.435]$ on extent and
$-0.564\,[-0.739, -0.392]$ on rate, both intervals excluding zero. The
extent contrast is precisely the seed-sensitive one noted on the
development set, and it held. Coverage is $0.88$--$0.91$ on all six
channels.

The secondary timing hypothesis was not supported. Timing calibration is
$-0.030$ $[-0.061, +0.006]$, the predicted sign but an interval that
includes zero, against a development basis of $-0.292$ whose interval had
excluded zero. It is accompanied by a small adverse timing recovery
difference whose interval excludes zero ($+0.003$ $[+0.000, +0.007]$,
three milliseconds against a 65-millisecond RMSE). The graph's value on
the timing channel therefore did not carry to the held-out test set. The development set's adverse loudness result did
\emph{not} replicate ($+0.015$, not significant), the delay channel
stayed graph-neutral as expected, and the ordering against the reference
targets is unchanged. The given variances remain the better-calibrated
noise model, with vibrato coverage $0.80/0.85$ against the learned
scale's $0.72/0.73$. By the decision rule fixed in advance, with the
three main hypotheses supported, the principal result of the
strings-and-winds setting holds: the unchanged graph prior extends to
real audio with calibrated uncertainty on a held-out test set. Every
number of the evaluation, all three systems and the per-entry evidence
are archived with the run logs.

\begin{table}[t]\centering
\caption{Strings-and-winds held-out results (the held-out test set
evaluated once, given-variance noise model, protocol fixed in advance),
paired per (track, seed) against the no-graph ablation ($95\%$ bootstrap
intervals, with $^{*}$ marking an interval that excludes zero and
$n = 78$--$80$ pairs per channel), with coverage at nominal $90\%$.
}
\label{tab:phase2-conf}
\footnotesize
\begin{tabular}{@{}lrrc@{}}
\toprule
channel & $\Delta$RMSE & $\Delta$NLL & cov@90 \\
\midrule
intonation $c$ (cents) & $-0.877^{*}$ & $-0.029$\phantom{$^{*}$} & 0.88 \\
extent $\log\gamma$ & $-0.425^{*}$ & $-2.990^{*}$ & 0.89 \\
rate $\log f^{\mathrm{vib}}$ & $-0.429^{*}$ & $-0.564^{*}$ & 0.89 \\
loudness $\ell$ & $-0.002$\phantom{$^{*}$} & $+0.015$\phantom{$^{*}$} & 0.90 \\
timing $\tau$ (s) & $+0.003^{*}$ & $-0.030$\phantom{$^{*}$} & 0.91 \\
delay $d^{\mathrm{vib}}$ (s) & $+0.006$\phantom{$^{*}$} & $+0.065$\phantom{$^{*}$} & 0.88 \\
\bottomrule
\end{tabular}
\end{table}

\figph{phase2_real_example.png}{1.0}{The six-channel bundle on one real development track
(violin), 30\% of notes hidden, four of the six channels shown. First
panel: intonation, with the estimator's per-note values and their own
uncertainties (dots), the posterior of the model with per-entry masks and
its 90\% predictive band, and the hidden notes' estimator values (open
circles). Second: vibrato extent, where squares mark notes whose extent
the estimator could not identify from the audio and the model supplies a
posterior anyway, informed by neighbours through the graph and by the
coupled channels. Third: the onset-anchored timing channel $\tau$, whose
error bars are the tempo line's leave-one-out predictive standard
deviations, the aligner error carried in the noise row (note the
wide-interval edge note at the right). Fourth: the vibrato onset delay
from the delayed-onset fit, where squares again mark estimator-missing
entries filled by the model. Its posterior may stray below zero on
unobserved stretches, since the channel is modelled as an unconstrained
Gaussian field while every observed delay is nonnegative by construction.
Loudness and the full per-channel numbers are in the tables. Loudness is
also the bundle's null result under the default noise
model.\label{fig:phase2-real}}

Figure~\ref{fig:arch-phase2} records the architectural change. Relative
to Figure~\ref{fig:arch-phase1}, the coregionalization grows from
$3 \times 3$ to $k \times k$ as channels join the bundle, re-learned
rather than carried over, and the observation noise $\Sigma_e$ becomes
substantial because the targets are estimates rather than symbolic MIDI
values. The prior's structure and the exact conjugate posterior are
untouched. The extension is scoped so that a negative result would cost
the piano contribution nothing: that contribution stands on its own, and
the extension's value is the demonstration that the same prior,
unchanged, extends to channels it was never designed for. That
demonstration is now made on real recordings and on a held-out test set
for the main hypotheses, while the secondary timing hypothesis was not
supported.

\figph{arch_phase2.pdf}{1.0}{The strings-and-winds system as run (development level,
\S\ref{sec:phase2-real}). Audio and the annotated onsets enter only to
derive the per-note entries $[c, \log\gamma,
\log f^{\mathrm{vib}}, \ell, \tau, d^{\mathrm{vib}}]$, each with
its own variance. One structural block changed relative to
Figure~\ref{fig:arch-phase1}, the likelihood, where $\Sigma_e$ is now
substantive and heteroscedastic, while the coregionalization is
re-learned at $k \times k$. The prior branches and the conjugate
posterior head are reused unchanged. The music-model kernel is absent
here, because the model is trained on piano, so the evidence branch
carries the score features alone.\label{fig:arch-phase2}}

\section{Two follow-up hypotheses on the held-out test set}
\label{sec:phase2-reuse}
Two hypotheses developed after the evaluation of \S\ref{sec:phase2-conf}
were tested on the same 13-piece held-out test set, which had therefore
already been used once: its targets and plain-graph fits were known when
the hypotheses were selected, while the hypothesis-specific quantities
had never been computed on it. Both hypotheses were developed on the
development set alone, with the primary measure, the decision rule, the
systems and the exact computation fixed before the test set was touched.
The inputs are the archived targets of the first evaluation, and the
masks, entries and seeds are those of Table~\ref{tab:phase2-conf}. This
second use of the held-out test set is weaker evidence than the single
use described in \S\ref{sec:hygiene}.\footnote{Reproduction records: the synthetic pilot of
\S\ref{sec:phase2-pilot} is recorded in
\texttt{docs/phase2\_synthetic\_pilot.md}, the two follow-up evaluations
in \texttt{results/corrnoise\_confirmation\_results.md} and
\texttt{results/phase2\_tonal\_confirmation\_results.md}, with their
evidence under \texttt{evidence/phase2\_reuse\_confirmation/}. The
correlated-noise computation below is the routine
\texttt{gp.noise\_corr} of Chapter~\ref{ch:future}.}

\subsection{Correlated timing noise}
The timing hypothesis of the main evaluation was not supported, and the
diagnosis pointed to the noise model: alignment error is correlated along
score time, which the diagonal noise row cannot represent. The revised
noise model places an AR(1) block in the timing channel's noise
covariance, $\Sigma_{\tau}(\rho)_{ij} = \sqrt{d_i d_j}\,\rho^{|i-j|}$ in
note order, with $d_i$ the given entry variances and $\rho$ chosen per
(track, seed) by the observed block's marginal likelihood on a fixed
grid. Held-out timing is predicted through the full covariance including
the noise cross-terms. Setting $\rho = 0$ recovers the published diagonal
path, which is the paired baseline, and the correlated-noise option of
Chapter~\ref{ch:future} is numerically identical to this computation. The
primary measure was the paired timing NLL, cluster-bootstrapped over the
40 unique tracks. The paired RMSE was a secondary measure that did not
enter the decision, and coverage and the share of (track, seed)
evaluations with $\rho > 0$ were descriptive.

The primary measure supported the hypothesis. The paired difference in
timing NLL is $-0.121$ with an interval of $[-0.156, -0.089]$, which
excludes zero, against a development basis of $-0.111$. The secondary
measure agrees: the paired RMSE difference is $-0.011$
$[-0.016, -0.007]$ s, from $0.058 \to 0.047$ s. Coverage stayed at
nominal, $0.908$ against the diagonal's $0.911$, and $\rho > 0$ on 79 of
80 (track, seed) evaluations, with median $0.45$. The risk stated in
advance, that this test set's timing channel is cleaner than the
development set's, showed as a smaller absolute recovery gain rather
than as an interval reaching zero. By the decision rule fixed in advance,
the correlated timing-noise row resolves the timing-calibration
hypothesis that the main evaluation had left unsupported, on this second
use of the held-out test set.

\subsection{The tonal metric}
The exploratory finding of \S\ref{sec:phase2-delta} and
Appendix~\ref{app:tonal}, that replacing the semitone metric of
\eqref{eq:adjacency} by the circle-of-fifths metric \eqref{eq:tonaladj}
improves held-out intonation, was tested with calibration as the primary
measure: the paired intonation NLL, tonal minus plain graph, under the
given variances, per (track, seed), with a 95\% bootstrap interval
required to exclude zero. Recovery, the paired intonation RMSE, was the
secondary measure and did not enter the decision, with its power stated
in advance at $0.52$--$0.63$ on this test set. The development timing
penalty was a descriptive prediction. All five arms ran together, and
the three plain arms are those of Table~\ref{tab:phase2-conf}.

The primary hypothesis was not supported. The paired intonation NLL
difference is $-0.023$ with an interval of $[-0.049, +0.001]$: the mean
lies on the predicted side but the interval reaches zero, against a
development basis of $-0.050$ $[-0.072, -0.029]$. The secondary measure
was supported, with a recovery difference of $-0.312$ cents
$[-0.599, -0.081]$ (development $-0.213$). The predicted timing penalty
did not appear under the given variances ($-0.002$ $[-0.011, +0.007]$).
By the decision rule fixed in advance, the thesis does not state that the
tonal metric improves intonation calibration. The metric stays
exploratory and Appendix~\ref{app:tonal} keeps its status. The three
plain arms reproduced Table~\ref{tab:phase2-conf} digit for digit, so
this evaluation left the earlier result untouched. Two descriptive
results are recorded beside the hypotheses: a small adverse loudness
calibration difference whose interval excludes zero ($+0.020$), and the
tonal graph improving on the no-graph ablation in intonation calibration
($-0.053$ $[-0.099, -0.010]$) where the plain graph had not
(\S\ref{sec:phase2-conf}: $-0.029$, not significant). The latter is
suggestive, was not a pre-specified hypothesis, and is not pursued further. Of
the two follow-up hypotheses, then, one was supported and one was not,
and the held-out test set has now been used twice.

% =====================================================================
